\section{Introduction}

A player owes \$1,400 and holds \$49 in cash. Its properties can cover the obligation, but some mortgages are unavailable until buildings have been sold. The player announces a liquidation budget, sells its development, mortgages eight deeds and pays, leaving the predicted \$14. Another player facing a smaller debt counts only the mortgages currently offered, overlooks collateral released by sales, and declares bankruptcy. These episodes pose a concrete question about LLM planning: does an agent anticipate how its actions change the resources available to it, and does its subsequent behavior realize that assessment?

The same question becomes strategic when the resources belong to someone else. Selling a deed can fund the seller's development while completing the buyer's monopoly. Withholding it preserves a blocker but forgoes cash. Financing a distressed rival can postpone a transfer to a stronger creditor. In these decisions, the relevance of an asset depends on ownership, available finance and what each participant can do next. A plausible account of those consequences is evidence of an expressed assessment. Its accuracy, realization and value require further evidence.

We study twenty complete MonopolyBench games played by four configured LLM systems. The records connect 6,989 decisions to delivered information, submitted actions and state changes. Our analysis covers all 920 negotiation responses and all 188 debt episodes requiring explicit financing. We use the complete populations to situate selected trajectories, and the trajectories to examine prospective statements alongside their outcomes. The scientific object is economic behavior in an interacting system: how agents describe opportunities, negotiate over them and execute local financial plans.

The clearest statement--action correspondence occurs when payment requires a change in the available collateral. Sixteen debts cannot be covered by current cash, all building salvage and the mortgages listed initially. Fourteen are paid after sales release additional mortgages. The initial text in thirteen successful cases states the additional finance or a budget including it; one considers it tentatively. Both failed cases omit released collateral in their calculations and immediately declare bankruptcy. The initial assessment and the subsequent actions together show how agents apply a financing rule to the state they face.

Bargaining has a different evidential structure. Its immediate choices are observable, but the value of accepting or rejecting a trade depends on future play. Opening swaps completing only their proposer receive 95/116 rejections, compared with 37/72 completing neither participant. The difference falls substantially when the proposer population changes. The cash-bid difference is negative only when pooled; within recipients it ranges from $-0.2$ to $+8.5$ percentage points. Separately, agents explicitly invoke rival construction, financing and creditor succession in their reasons and messages. We report these as complementary observations rather than treating the reasons as an explanation of the aggregate contrast.

Eight contextual incidents show the range of conduct behind the counts. They include refusing one's own monopoly while citing the rival's opportunity, pricing a completing deed, bluffing about a limit, correcting a false pitch, selectively describing financing, seeking collective financial denial, withdrawing an assurance and carrying out a payment plan. Quotations identify their channels and are connected to delivered messages and subsequent events.

Our contributions are a complete account of the encountered negotiation and financing decisions; an analysis of financing dependencies linking prospective statements to executed payment; and contextual evidence of how strategic communication operates over persistent economic stakes. The design supports those observational claims. It does not provide an independent value reference for trades or a comparison of underlying model intelligence. The five seed blocks, response conditions and retrospective analysis choices are made explicit so that the scope of each finding is visible.

The paper first defines the observed systems and measurement populations, then presents financial planning, bargaining and conduct. Its discussion connects these findings to the motivating questions about strategic reasoning, opponent modeling, negotiation, memory, adaptation, economic judgment, collective behavior and evaluation validity. The appendix retains the full question map, episode and offer tables, prompts and artifact entry points.
